\section{Lyndon's homomorphic-image theorem: proofs}
\label{sec:annex-p03-n-a-011-02}

The imported polarity interpolation theorem and its stronger constructive
Main Theorem have their stated source boundary in
\cref{thm:p03-lyndon-interpolation-import,thm:p03-lyndon-constructive-import}.
The proofs here reconstruct the preservation article's arguments, with
explicit assignment, quotient, and compactness steps
\cite[pp.~145--153]{Lyndon1959}.

\begin{annexitem}{annex:p3.n-a-011.2.increasing}{thm:p03-lyndon-increasing}{Polarity, interpolation, and increasing theories}
\begin{proof}
First put a formula in negation normal form. Pushing negation through
connectives and quantifiers, deleting double negations, and leaving
\(0,1\) as constants preserves the polarity of each surviving relation
occurrence. A \(Q\)-positive literal is either an atom, or a negated atom
whose relation lies outside \(Q\). The first is preserved when the
selected relations increase, and the second is unchanged. Conjunction
and disjunction preserve this implication. The same assignments witness
existential quantifiers before and after enlargement; universal
quantifiers use the same domain and all the same assignments. Induction
therefore proves the first assertion, including formulas with free variables.

Fix \(D\in\Delta\). Compactness gives finite subsets of the premises
entailing \(D'\). Choose a single finite subset of \(\Sigma\) large
enough to cover both the unprimed and primed occurrences, and let \(S\)
be its conjunction. Let \(C\) be the conjunction of the selected members
of \(\Gamma\), and \(J\) the selected inclusion sentences. Then
\[
 \{S,S',C,J\}\entails D',\qquad
 \{S\land C\}\entails ((J\land S')\to D').
\]
The right-hand sentence has only positive occurrences of members of
\(Q\): their negative occurrences in the antecedent \(J\) are reversed
by the outer implication. The symbols \(Q'\) do not occur on the left.
The imported interpolation theorem supplies \(P_D\), positive in \(Q\)
and containing no \(Q'\), between these two sentences. Thus
\(\Sigma\cup\Gamma\entails P_D\). Replace each \(q'\) by \(q\)
in the implication from \(P_D\). The resulting \(J\) is valid, so
\(\{P_D,S\}\entails D\). Set \(\Pi=\{P_D:D\in\Delta\}\).
No positivity assumption on \(\Gamma\) is needed: it is the polarity
on the right side of the interpolation application that controls \(P_D\).

If \(\Gamma\) is increasing, apply this assertion with empty
\(\Sigma\) and \(\Delta=\Gamma\), obtaining equivalent
\(\Gamma,\Pi\). Conversely, if \(\Gamma\) is equivalent to a
\(Q\)-positive theory \(\Pi\), simultaneous symbol replacement gives
equivalent \(\Gamma',\Pi'\). The first assertion gives
\(\Pi\cup I_Q\entails\Pi'\), hence
\(\Gamma\cup I_Q\entails\Gamma'\).
\end{proof}
\end{annexitem}

\begin{annexitem}{annex:p3.n-a-011.2.maps}{thm:p03-lyndon-map-calculus}{All map and quotient steps}
\begin{proof}
Term induction gives
\(h(t^A[s])=t^B[h\circ s]\), including zero-ary operations.
For an atom use preservation of its relation; for a negated atom outside
\(Q\), use reflection. The Boolean cases follow directly. An existential
witness \(a\) maps to the witness \(h(a)\). For a universal formula
and any \(b\in B\), choose \(a\in A\) with \(h(a)=b\); apply the
induction hypothesis to the assignment changed to \(a\). This is the
precise use of surjectivity. For a \(\varnothing\)-map the same argument,
now with atomic equivalence and Boolean negation, proves equivalence for
all formulas; in the existential reverse direction choose a preimage of
the witness, and in the universal reverse direction evaluate at \(h(a)\).

Let \(k(a,a')\) mean \(h(a)=h(a')\). It is an equivalence relation.
If corresponding inputs are \(k\)-equivalent, operation preservation
makes their outputs \(k\)-equivalent. Atomic reflection makes each
relation's truth unchanged by replacing inputs by \(k\)-equivalent
ones. These are exactly the congruence conditions. For an arbitrary such
congruence, define quotient operations on classes by applying the operation
to representatives, and quotient relations by the relation on
representatives. The stated conditions prove representative independence.
The canonical quotient surjection is a \(\varnothing\)-map. For the
kernel of \(h\), send \([a]_k\) to \(h(a)\). This is well defined,
injective, onto, and preserves and reflects the interpretations, giving
the required factorization.
\end{proof}
\end{annexitem}

\begin{annexitem}{annex:p3.n-a-011.2.diagram}{thm:p03-lyndon-diagram}{Quotienting the expanded complete diagram}
\begin{proof}
Let \(B_A\) be the given expansion. Every sentence saying \(e_A\) is
an equivalence relation, is substitutive for every operation, and preserves
every relation under replacement of coordinates is true in \(A_A\).
All these sentences lie in \(A_A^*\); consequently \(k=e_A^{B_A}\)
is a congruence. The quotient \(C_A=B_A/k\) still satisfies the whole
diagram, by the formula equivalence for a \(\varnothing\)-map, and
\(e_A^{C_A}\) is actual identity.

Define \(j(a)=w_a^{C_A}\). For distinct \(a,a'\), the diagram contains
\(\neg e_A(w_a,w_{a'})\), so \(j\) is injective. If
\(w^A(\bar a)=a'\), the diagram contains
\(e_A(w(\bar w_a),w_{a'})\), proving preservation of operations.
For every relation \(r\) and tuple \(\bar a\), exactly its true
atomic statement or its negation belongs to the complete diagram; thus
\(r^A(\bar a)\) iff \(r^{C_A}(j\bar a)\). This identifies \(A_A\)
with the named substructure of \(C_A\).

For an arbitrary \(L\)-formula \(F(\bar x)\), substitute the indicated
constants for its free variables. Both truth and falsity of the resulting
sentence in \(A_A\) are recorded in its complete theory. Since \(C_A\)
satisfies that theory, \(F\) has the same truth at \(\bar a\) in
\(A\) as at \(j\bar a\) in the reduct \(C\). Hence \(j\) is an
elementary embedding. Relabel \(j(a)\) by \(a\), using fresh names
for other elements. The relabeled reduct is an elementary extension of
\(A\), and the quotient map from \(B\) followed by this relabeling
is the required \(\varnothing\)-map. Its direction is from \(B\)
to the extension, not the reverse.
\end{proof}
\end{annexitem}

\begin{annexitem}{annex:p3.n-a-011.2.qimage}{thm:p03-lyndon-qimage}{Interpolation, consistency, and the common-domain construction}
\begin{proof}
Put \(\Gamma=K^*\), \(\Delta=A_A^*\), and leave the new constants
and \(e_A\) unchanged under priming. Suppose
\(\Gamma\cup I_Q\cup\Delta'\) were inconsistent. Compactness gives
a finite conjunction \(D\) from \(\Delta\) such that
\(\Gamma\cup I_Q\entails\neg D'\). The preceding interpolation
argument, with empty \(\Sigma\), supplies a \(Q\)-positive sentence
\(P\) with
\[
 \Gamma\entails P,\qquad \{P\}\entails\neg D.
\]
The fresh relation \(e_A\) is absent from the first side, hence absent
from \(P\) by the polarity condition. If named constants occur in
\(P\), replace their finitely many occurrences by fresh variables and
universally close. Call the result \(P_0\). Since the constants are
absent from \(\Gamma\), arbitrary interpretations of them show
\(\Gamma\entails P_0\); and \(P_0\entails P\) with their original
interpretations. Thus \(P_0\) is a \(Q\)-positive \(L\)-sentence
entailed by \(\Gamma\). The hypothesis gives \(A\entails P_0\),
hence \(A_A\entails P\), contradicting \(D\in A_A^*\).
This makes explicit the removal of diagram parameters in the source proof.

Take a structure \(C\) satisfying the now-consistent combined theory.
On its domain define \(G\) by the unprimed \(L\)-interpretations,
and \(B_A\) by using \(q'^C\) for \(q\in Q\), keeping operations,
other relations, constants, and \(e_A\) unchanged. Then \(G\entails
\Gamma\), so \(G\in K\), and \(B_A\entails\Delta\).
The identity of their common domain is a \(Q\)-map from \(G\) onto
the \(L\)-reduct \(B\): selected relations increase by \(I_Q\),
others agree. Proposition~3 gives a \(\varnothing\)-map from \(B\)
onto an elementary extension \(A'\) of \(A\). Composition gives a
\(Q\)-map \(G\to A'\), including reflection for relations outside
\(Q\).

Conversely let \(G\in K\) map onto an elementary extension \(A'\)
by a \(Q\)-map. The map lemma preserves each \(Q\)-positive sentence
of \(K^*\) to \(A'\); elementary equivalence then transfers it to
\(A\). To see the source's relation-enlargement reasoning for a map
whose domains differ, pull the target relations back along the surjection
to the source domain. This increases the selected relations, leaves the
other ones fixed, and does not alter operations. The pullback structure
maps to the target by a \(\varnothing\)-map. Proposition~1 applies
on the common domain, and the map lemma transfers the resulting truth
to the target. Identifying the domains without this step would not
justify a noninjective map.
\end{proof}
\end{annexitem}

\begin{annexitem}{annex:p3.n-a-011.2.qclosure}{thm:p03-lyndon-qclosure}{Axiomatization by positive consequences}
\begin{proof}
Suppose \(K\) is elementary and closed under \(Q\)-images. Let
\(\Pi\) be its \(Q\)-positive common sentence theory. Certainly
\(K\subseteq\Pi^*\). For \(A\in\Pi^*\), Theorem~4 gives an
elementary extension \(A'\) which is a \(Q\)-image of some \(G\in K\).
Closure gives \(A'\in K\); elementary equivalence gives \(A\in K\)
since \(K\) is axiomatized by sentences. Thus \(K=\Pi^*\).
Conversely preservation of every \(Q\)-positive axiom gives closure.
Apply this equivalence to \(K=\Gamma^*\) for the theory formulation.
\end{proof}
\end{annexitem}

\begin{annexitem}{annex:p3.n-a-011.2.main}{thm:p03-lyndon-main}{The three quotient cases with genuine identity}
\begin{proof}
Let \(Q_1=R\), \(Q_2=\{e\}\), and \(Q_3=R\setminus\{e\}\).
Any map in the respective row is a \(Q_i\)-map of auxiliary structures,
so the forward preservation implication follows from Theorem~4.

For the converse apply that theorem to the auxiliary elementary closure
\(K^{**}\), whose common theory is still \(K^*\). It gives an
auxiliary \(C\in K^{**}\) and a \(Q_i\)-map
\(h:C\to B\), where \(B\) is an elementary extension of \(A\).
All congruence axioms for \(e\) hold throughout \(K\), hence in \(C\).
They hold in \(B\) as well. Although \(e^A\) is identity,
\(e^B\) may still identify new elements: ``identity'' here is not an
auxiliary-language sentence. Quotient both structures by these congruences.
The quotient \(\bar C=C/e^C\) is a relational system with the same
sentence theory as \(C\), by the \(\varnothing\)-map lemma. Hence
\(\bar C\in K\). Similarly \(\bar B=B/e^B\) is a relational
system. Its named copy of \(A\) remains injective since the restriction
of \(e^B\) to \(A\) is identity. Quotient formula equivalence and
\(A\preccurlyeq B\) prove that \(\bar B\), after relabeling, is
an elementary extension of \(A\).

For all three cases the induced map
\(\bar h([c]_{e^C})=[h(c)]_{e^B}\) is well defined and onto, and
preserves operations and relations. For \(Q_1\) this is the required
onto homomorphism. For \(Q_2\), reflection of all nonequality relations
shows that \(\bar h\) is a projection: its kernel is substitutive
for operations and those relations, and the quotient by the kernel is
isomorphic to \(\bar B\). For \(Q_3\), reflection of \(e\) says
\(e^B(h(c),h(c'))\) iff \(e^C(c,c')\), so \(\bar h\) is
injective. Thus it is exactly an enlargement. This induced-map argument
also explains the source's two-stage quotient diagram without assuming
that the original \(h\) was injective.

Finally, for any onto homomorphism between relational systems, replace
each nonequality relation of its source by the full inverse image of
the target relation. This is an enlargement; the map now reflects those
relations, and its kernel supplies the subsequent projection.
\end{proof}
\end{annexitem}

\begin{annexitem}{annex:p3.n-a-011.2.positive}{thm:p03-lyndon-positive-corollaries}{All directions of the three positive corollaries}
\begin{proof}
Let \(\Pi\) be the positive sentences true in \(K\). If a positive
\(P\) true in \(K\) entails \(S\), map preservation gives \(P\),
hence \(S\), in every image. Conversely suppose \(S\) holds in
every image of \(K\). For \(A\entails\Pi\), Theorem~5 gives an
elementary extension \(A'\) that is such an image. Then
\(A'\entails S\), so \(A\entails S\). Hence
\(\Pi\entails S\), with consequence among relational systems.
Compactness in the ordinary equality semantics gives finitely many members of
\(\Pi\) entailing \(S\). Their conjunction is the required positive
\(P\). The empty conjunction is \(1\).

Suppose \(K\) is image-closed and \(K=\Gamma^*\) among relational
systems. Each \(S\in\Gamma\) holds in its images, so the preceding
argument makes \(S\) a consequence of some \(P_S\in\Pi\).
Therefore \(\Pi^*\subseteq\Gamma^*=K\); the other inclusion is
automatic. Conversely positive axioms are preserved, giving image closure.

For a single image-preserved \(S\), take \(K=\{S\}^*\). The first
corollary gives positive \(P\) true throughout \(K\) and entailing
\(S\). Thus \(S\) entails \(P\) and \(P\) entails \(S\).
The opposite implication follows directly from positive preservation.
This also covers an inconsistent \(S\), whose positive equivalent is
\(0\); no consistency hypothesis is hidden in the sentence result.
\end{proof}
\end{annexitem}

\begin{annexitem}{annex:p3.n-a-011.2.two-extensions}{thm:p03-lyndon-two-extensions}{Naming the second structure}
\begin{proof}
The backward direction follows from positive preservation and elementary
equivalence. For the forward direction expand the language by constants
for all elements of \(B\), and let \(K_B\) be the elementary class
defined by the complete elementary diagram of \(B\), with genuine
identity. Any positive sentence \(P(\bar c)\) true in \(K_B\) is
true in the named \(B\). Its existential closure in the finitely many
displayed constants is a positive original-language sentence true in
\(B\), and therefore in \(A\).

The elementary diagram of \(A\), together with all these positive
consequences of \(K_B\), is finitely satisfiable: conjoin the finitely
many positive sentences concerned and existentially close their finitely
many \(B\)-constants. A witness in \(A\) interprets those constants,
while the \(A\)-diagram is satisfied by its own names. Compactness
gives an elementary extension \(A_0\) of \(A\) with an expansion
satisfying the full positive theory of \(K_B\). Theorem~5 in this
expanded language now gives an elementary extension of that expansion
as an image of a member \(B'\in K_B\). Forget the constants to obtain
\(A'\) and the onto homomorphism \(B'\to A'\). The diagram makes
the named copy of \(B\) elementary in \(B'\); relabel it if necessary.
Transitivity of elementary extension gives \(A\preccurlyeq A'\).
\end{proof}
\end{annexitem}

\begin{annexitem}{annex:p3.n-a-011.2.relative}{thm:p03-lyndon-relative}{Relative preservation without assuming the ambient class is image-closed}
\begin{proof}
Theorem~5 already gives the main equivalence: \(K\) is elementary,
and the elementary extension of \(A\in H\) remains in \(H\).
For the relative first corollary, suppose \(S\) holds in all images
of \(K\) that lie in \(H\). An \(H\)-system satisfying the positive
theory \(\Pi\) of \(K\) has an elementary extension that is such
an image. Consequently \(H\cup\Pi\entails S\), where \(H\)
denotes any sentence axiomatization of the ambient class. Compactness
gives a finite positive conjunction \(P\in\Pi\) with
\(H\cup\{P\}\entails S\). Conversely \(P\) is preserved into
every image, and the \(H\)-hypothesis supplies the relative implication.

For a relatively image-closed elementary \(K\subseteq H\), apply
this consequence assertion to each axiom of \(K\). Its positive
consequences axiomatize \(K\) within \(H\). Conversely such a
relative positive axiomatization is preserved whenever both endpoints
are in \(H\). For \(K=\{A\in H:A\entails S\}\), the same finite
conjunction argument gives \(H\)-equivalence to one positive sentence.
The converse is positive preservation. This proves all relative
directions even when \(H\) is not itself closed under images.
\end{proof}
\end{annexitem}

\begin{annexitem}{annex:p3.n-a-011.2.algebra}{thm:p03-lyndon-algebra-examples}{Concrete multiplicative-system witnesses}
\begin{proof}
For a surjection \(h\) from a commutative system, choose preimages
\(a,b\) of arbitrary \(u,v\). Then
\(uv=h(a)h(b)=h(ab)=h(ba)=h(b)h(a)=vu\).
For cancellation take the cancellative additive semigroup
\((\mathbb N,+)\) and its two-element image with operation
\(u\mathbin{\vee}v=\max(u,v)\), under the map sending zero to zero
and every positive integer to one. The map is onto and preserves the
operation, since \(m+n\) is positive iff at least one input is positive.
In the image \(1\vee0=1\vee1\), although \(0\ne1\). Thus left
cancellation fails in an image of a cancellative system. The preservation
theorem rules out any positive equivalent, not merely this particular
implication-shaped presentation of cancellation.
\end{proof}
\end{annexitem}

\begin{annexitem}{annex:p3.n-a-011.2.nonelementary}{thm:p03-lyndon-nonelementary}{Finite pure sets and their common theory}
\begin{proof}
For a pure identity formula of quantifier rank at most \(m\), compare
assignments with the same equality pattern in two sets, each having at
least \(m\) as-yet-unassigned elements. Atomic identity has the same
truth. Boolean connectives preserve agreement. For a quantifier, a
witness either is an already assigned element, in which case match it,
or is new, in which case choose an unused element in the other set.
The remaining supply and quantifier rank both decrease by one. Induction
proves agreement for the formula. Thus every sentence true in all finite
nonempty sets is true in every infinite set: compare the infinite set to
a finite set large enough for its quantifier rank. Its truth in finite
sets is already assumed, so it is valid in the pure identity semantics.
In particular every set satisfies their common positive theory.

An elementary extension of an infinite set contains that set, so is
infinite. No surjection from a finite set can have this extension as
codomain. The class of finite sets is image-closed but is not elementary,
and fails precisely the conclusion that requires elementarity.
\end{proof}
\end{annexitem}

\begin{annexitem}{annex:p3.n-a-011.2.chain}{thm:p03-lyndon-image-class-boundary}{The chain obstruction and the source's matrix argument}
\begin{proof}
If \(B\entails S\), choose its root \(b_0\) and first neighbor
\(b_1\). Repeatedly choose \(b_{i+1}\) from the quantified extension
condition for \(b_i\). Dependent choice gives
\(r(b_0,b_i)\) for \(i\geq1\) and
\(r(b_i,b_{i+1})\) for \(i\geq1\). The terms need not be distinct.
An onto homomorphism transports this infinite rooted chain to its image.

Conversely, if \(A\) contains such a sequence \((a_i)_{i\geq0}\),
adjoin to a disjoint copy of its domain distinct elements \(b_i\).
Give the new structure exactly the edges \((b_0,b_i)\), \(i\geq1\),
and \((b_i,b_{i+1})\), \(i\geq1\), and no edges on its copied
old domain. It satisfies \(S\). Map the copied domain identically to
\(A\) and \(b_i\) to \(a_i\). This is an onto homomorphism.
Hence \(H(K)\) is exactly the class containing a rooted infinite chain.

Let \(A\entails S_n\) for all \(n\). Add new constants \(d_i\)
to its elementary diagram and all the rooted-chain edge requirements.
Every finite subset holds in \(A\), interpreting the finitely many
new constants along a sufficiently long \(S_n\)-chain. Compactness
gives an elementary extension \(A'\) containing an infinite chain.
Thus \(A'\in H(K)\), and every sentence common to \(H(K)\) holds
in \(A\). Conversely all the \(S_n\) hold throughout \(H(K)\),
and therefore throughout its elementary closure. This proves the
displayed elementary-closure identity. If \(S\entails P\) for
positive \(P\), preservation gives \(P\) throughout \(H(K)\).
The identity implies \(\{S_n:n\geq1\}\entails P\). Compactness
gives a finite subset, and the longest chain sentence in it entails all
the others; hence \(S_n\entails P\) for some \(n\).

For completeness, the source's p.~152 derivation of this last fact uses
the constructive interpolation dependency rather than this diagram
argument. The Skolem matrix of \(S\), with fresh constants \(s_0,s_1\)
and a fresh unary operation \(s\), is
\[
 r(s_0,s_1)\land
 (r(s_0,z)\to(r(s_0,s(z))\land r(z,s(z)))).
\]
Apply the imported Main Theorem to \(S\) and \(\neg P\), then its
literal deletion construction. The second matrix is negative, so the
retained first matrix is positive. It yields an intermediate positive
sentence \(U\) entailing \(P\), whose truth follows from a universal
positive matrix \(M\). In the source's reconstruction, \(M\) is a
propositional consequence of a set \(\Sigma\) of substitution instances
of the displayed Skolem matrix.

Put \(t_0=s_0,t_1=s_1,t_{i+1}=s(t_i)\) for \(i\geq1\).
Define an atomic Boolean substitution \(\chi\) which keeps precisely
the atoms \(r(t_0,t_i)\), \(i\geq1\), and
\(r(t_i,t_{i+1})\), \(i\geq1\), and replaces the other atoms by
\(0\). It is a substitution on propositional generators, not a
substitution for individual variables. Positivity gives
\(\chi M\entails M\). Since propositional consequence is preserved
under substitution, \(\chi\Sigma\entails\chi M\), hence
\(\chi\Sigma\entails M\). Each retained instance contributes
\(r(t_0,t_1)\) and an implication that extends the retained rooted
path; all other implications have false antecedent. Thus the substituted
instances are equivalent to the full set of rooted-chain atoms.
Finite propositional compactness reduces the implication to a finite
initial rooted chain, giving the same \(S_n\entails U\entails P\)
under the imported duplication/specialization reconstruction. This is
the precise constructive dependency of the original argument; the
preceding diagram proof separately verifies the stated conclusion
without promoting that external construction to a proof supplied here.

Finally on \(\mathbb N\) let \(r(i,j)\) hold iff \(j<i\).
For every \(n\), \(n,n-1,\ldots,0\) witnesses \(S_n\).
There is no infinite strictly decreasing sequence of natural numbers,
so there is no rooted infinite chain and the structure is not in \(H(K)\).
It is in \(H(K)^{**}\), proving \(H(K)\) is not elementary and
showing why Theorem~5 must refer to an elementary extension of \(A\).
\end{proof}
\end{annexitem}

\begin{annexitem}{annex:p3.n-a-011.2.undecidable}{thm:p03-lyndon-undecidable}{The recursive one-element test and reduction}
\begin{proof}
Write \(C\) for a finite conjunction of axioms of \(T\), and let
\(\Sigma\) consist of the sentences true in every one-element
\(T\)-system. Membership of \(\Sigma\) is decidable. A sentence
\(C\to D\) mentions only finitely many relation symbols; in a
one-element domain each operation is uniquely determined and each
mentioned relation has just one truth value. Enumerate the finitely
many truth assignments to these relations, fixing identity to true,
and evaluate \(C\to D\). Quantifiers range over the single element.

Every theorem of \(T\) belongs to \(\Sigma\). If theoremhood were
decidable on \(\Sigma\), the recursive membership test, followed by
that restricted procedure, would decide it on all sentences: outside
\(\Sigma\) the answer is automatically no. Thus theoremhood remains
undecidable on \(\Sigma\).

For \(D\in\Sigma\), if \(T\entails D\), then \(\neg D\) is
\(T\)-equivalent to the positive false sentence \(0\). If
\(T\not\entails D\), there is a \(T\)-system \(A\) with
\(A\entails\neg D\). By hypothesis it has a one-element onto image
\(B\) satisfying \(T\). Since \(D\in\Sigma\),
\(B\entails D\). Therefore \(\neg D\) is not relatively preserved
and, by Theorem~5', has no \(T\)-equivalent positive sentence. We have
proved that \(\neg D\) has such an equivalent iff \(D\) is a theorem,
for \(D\in\Sigma\). An algorithm deciding relative positive
equivalence would decide this impossible restricted theoremhood problem.

For groups the unique map to the trivial group preserves the group
operations and identity, so the image hypothesis holds. In the pure
predicate example, interpret all nonequality predicates true on a
singleton; the constant map preserves every relation. The undecidability
of the selected theory in each example remains the separately cited
premise identified on source p.~153, not a consequence of the image
construction itself.
\end{proof}
\end{annexitem}

\begin{annexitem}{annex:p3.n-a-011.2.boundaries}{thm:p03-lyndon-lookup-boundaries}{Onto maps, arbitrary maps, falsum, and tuples}
\begin{proof}
For existential-positive formulas, the term, atom, Boolean, and witness
steps of the map lemma do not require preimages of arbitrary target
elements. They therefore work for homomorphisms that are not onto.
For universal formulas the preimage step is essential. In the unary
example in the body, the inclusion preserves all true source atoms and
equality, but the new target element falsifies \(\forall x\,P(x)\).
This is an explicit separation, not a claim that the original paper
proves the arbitrary-homomorphism characterization.

For a formula \(F(\bar x)\), replace its free variables by fresh
constants \(\bar c\). A homomorphism preserves their interpretations
exactly when it transports the chosen tuple. Sentence preservation in
all such expansions gives a positive \(P(\bar c)\) equivalent to
\(F(\bar c)\). Replace these constants back by distinct free variables,
renaming bound variables first. Arbitrary interpretations of the
constants prove the formula equivalence, and positivity is unchanged.
The reverse implication is the positive formula induction.

If \(0\) is omitted from the positive pure identity fragment, every
atom is true in a singleton and induction makes every positive sentence
true there. Thus no such sentence can be equivalent to an inconsistent
sentence. Including \(0\) supplies exactly that case. Finally the
compactness constructions produce arbitrary structures and extensions;
neither their hypotheses nor their conclusions impose finiteness. A
finite-only theorem cannot discharge these all-structure statements.
\end{proof}
\end{annexitem}
